% ============================================================
% Verwendete KI-Tools (Handbuch Kap. 7.4.2, S. 92)
% "Die Verwendung von KI-Tools muss in einem zusätzlichen Abschnitt am
% Ende des normalen Quellenverzeichnisses nachgewiesen werden [...]:
% Name des Tools, Datum der Verwendung, eingegebene Aufforderungen
% (Prompts)." Zusätzlich ist ein Reflexionsformular einzureichen (siehe
% Hinweis in selbstaendigkeitserklaerung.tex).
%
% Diese Angaben nach Bedarf ergänzen, falls im Verlauf der Arbeit weitere
% KI-Tools oder Prompts hinzukommen.
% ============================================================

\chapter*{Verwendete KI-Tools}
\addcontentsline{toc}{chapter}{Verwendete KI-Tools}

Das vorliegende LaTeX-Dokumentgerüst (Formatierung gemäss Tab.\,14,
Kapitelstruktur gemäss Kap.\,7.1.2, Zitier- und
Bibliografie-Vorlagen gemäss Kap.\,7.5) wurde mit Unterstützung des
KI-Tools Claude erstellt. Der wissenschaftliche Inhalt der Arbeit
(Fragestellung, Durchführung, Auswertung, Diskussion) wird eigenständig
von den Autoren verfasst; der Beispieltext in der Einleitung wurde als
Platzhalter mit demselben Tool entworfen und ist vor Abgabe durch
eigene Formulierungen zu ersetzen.

\begin{itemize}
  \item \textbf{Name des Tools:} Claude Code (Claude Sonnet 5), Anthropic
  \item \textbf{Datum der Verwendung:} 10.\,September 2026
  \item \textbf{Eingegebene Aufforderungen (Auszug):}
    \begin{enumerate}
      \item[a)] \enquote{Baue mir, gemäss den Vorgaben im Handbuch, ein
        LaTeX-Grundgerüst für die Probearbeit zum Thema Bierbrauen
        (naturwissenschaftliche Arbeit, Kurzbelege in Klammern).}
      \item[b)] \enquote{Ergänze die Aufsätze/Zeitschriften-Treiber im
        Quellenverzeichnis nach demselben Muster.}
      \item[c)] \enquote{Schreibe einen kurzen Beispieltext in die
        Einleitung und füge eine Referenz auf das Buch
        \emph{bruecklmeier2022} ein.}
    \end{enumerate}
\end{itemize}
